\subsection{Study design and reporting}
\label{sec:design}
This is a retrospective, cross-sectional survey study of competitive athletics athletes affiliated with the Federazione Italiana di Atletica Leggera (FIDAL). Reporting follows the STROBE statement and its athletics-relevant STROBE Sports Injury and Illness Surveillance (STROBE-SIIS) extension \citep{vonelm2007strobe,bahr2020strobesiis}; the completed checklist, with the manuscript line number supporting each item, is given in Appendix~\ref{sec:strobe-checklist}.

\subsection{Study registration}
\label{sec:registration}
This study was not registered in a public trial or study registry (for example, ClinicalTrials.gov or an equivalent observational-study registry). It was, however, prospectively specified through a complete study protocol, submitted in Italian and reviewed and approved by the Comitato Etico of the University of Turin (Section~\ref{sec:setting}) before survey deployment. The protocol specified the study population, the two-phase recruitment and poststratification strategy, the exposure and outcome definitions, and the primary hypothesis (Section~\ref{sec:introduction}) in advance of data collection and of any outcome analysis. This ethics-committee-reviewed protocol is the prospective record of the study design; it is not, however, a substitute for public registration, and the absence of the latter is noted as a limit on independent, external verification of pre-specification (Section~\ref{sec:discussion-limitations}).

\subsection{Setting and ethics}
\label{sec:setting}
Data collection used a structured, self-administered online survey covering demographic, anthropometric, training, footwear, and injury-history items for the current competitive season, separately for the preparation and competition periods. Study data were collected and managed using REDCap electronic data capture tools hosted at the Department of Clinical and Biological Sciences, University of Turin \citep{harris2009redcap,harris2019redcap}. REDCap (Research Electronic Data Capture) is a secure, web-based software platform designed to support data capture for research studies, providing (1) an intuitive interface for validated data capture, (2) audit trails for tracking data manipulation and export procedures, (3) automated export procedures for seamless data downloads to common statistical packages, and (4) procedures for data integration and interoperability with external sources.

The study was approved by the Comitato Etico of the University of Turin (protocol n. 0453001, approved 23 June 2026). Participation required electronic informed consent; no financial incentive was offered. Individual survey responses were stored under a study identifier disconnected from athlete names at the point of analysis; only aggregate, de-identified data and model outputs are reported here, in line with the data-privacy provisions of the Declaration of Helsinki.

\subsection{Participants and two-phase recruitment}
\label{sec:participants}
Eligible athletes were FIDAL-affiliated competitive athletes in the Under 20, Under 23, or Senior (\emph{assoluta}) categories, across seven athletics disciplines with adequate cell support in the sample (sprint, middle distance, long distance, hurdles, horizontal jumps, pole vault, combined events; four further disciplines had zero respondents and were dropped from discipline-level adjustment; Appendix Table~\ref{tab:discipline-support}).

Recruitment proceeded in two phases to reduce the risk that a purely voluntary sample would over-represent already engaged or higher-exposure athletes. In Phase 1, the survey was distributed openly through FIDAL-affiliated clubs and athlete networks; enrollment was entirely voluntary. National population totals by age category, sex, and discipline group were estimated from FIDAL competition-registry history (Appendix Table~\ref{tab:population-frame-history}). Study size was planned for precision of a population proportion rather than power for a specific exposure association. For each candidate sample size $n$, the worst-case count $x=\lfloor n/2\rfloor$ was combined with a Beta(1,1) prior, giving a Beta$(1+x,1+n-x)$ posterior; the smallest $n$ whose equal-tailed 95\% credible interval had half-width no greater than 0.06 was selected. This rule gave $n=264$ from the estimated eligible frame of 1328 athletes (19.9\%); targets were allocated proportionally across the 42 age-by-sex-by-discipline strata using largest-remainder rounding (Appendix Table~\ref{tab:sample-size}). The calculation established a minimum precision and population-coverage target, not representativeness by itself. After Phase 1 closed, strata that had not reached their target count were identified, and Phase 2 undertook targeted recall--direct outreach to athletes and clubs within under-represented age-by-sex-by-discipline strata--until every stratum met or approached its target or the recruitment window closed. Detailed recruitment-monitoring cells are retained privately to avoid disclosing small-cell participation patterns. This two-phase design does not produce a probability sample with known, non-zero, equal selection probabilities in every stratum; population coverage was supported by targeted recruitment, and residual over- and under-representation was corrected analytically (Section~\ref{sec:weighting}) rather than assumed away.

Of all survey respondents, 355 provided consent, met calculated eligibility, and completed the survey (\texttt{survey\_atleta\_it\_complete = 2}). One athlete was excluded for incomplete injury-timing data, and two further athletes were excluded because a period (preparation or competition) had a zero total eligible training-frequency score, precluding a carbon-share denominator. The final balanced analysis cohort comprised 352 athletes contributing 704 athlete-period records (one preparation and one competition record per athlete; Table~\ref{tab:sample-flow}).

\subsection{Exposure: carbon-plate insole use}
\label{sec:exposure}
For each period (preparation, competition) and each of five training domains--long endurance, middle endurance, lactate, maximum velocity, and technique--athletes reported an ordinal weekly training-type frequency (0, 1, 2, 3, 4, or 5-or-more sessions per week, with 5-or-more coded as 5) separately for sessions performed in carbon-plated spikes and sessions performed without them. \emph{Carbon share} for a given period was defined as
\[
\text{carbon share} = \frac{\sum \text{carbon-frequency items}}{\sum \text{carbon-frequency items} + \sum \text{non-carbon-frequency items}},
\]
summed across the five domains. The corresponding \emph{absolute carbon-frequency score} was the numerator alone, that is, the sum of the five carbon-frequency items. It is an ordinal frequency score rather than an exact count of unique sessions because responses were top-coded at five and a session could contribute to more than one reported training domain. Gym/strength conditioning (\emph{palestra}) has no carbon-eligible equivalent and was retained only as a non-carbon adjustment covariate, not in the carbon-share denominator. Carbon share is therefore a composition of reported ordinal training-type frequency, not a directly measured training-time proportion or shoe-wear duration; this distinction is carried through to the interpretation of all exposure estimands. Exposure was decomposed into two components: \emph{adoption} (any carbon use versus none, \texttt{carbon\_frequency\_sum > 0}) and, among users, continuous \emph{dose}, expressed primarily as carbon share and additionally as the absolute carbon-frequency score.

\subsection{Outcome: period-specific injury}
\label{sec:outcome}
The outcome was self-reported injury attributable to athletics training or competition participation during the current season, following the athletics-specific consensus definition \citep{bahr2020strobesiis,timpka2014consensus}: any physical complaint or observable tissue damage resulting from the transfer of energy experienced by an athlete during participation, regardless of medical attention sought or time-loss consequence. Athletes reported the number of current-season injuries (0, 1, or 2) and, for each, whether onset occurred in the preparation period, the competition period, or both (coded 1, 2, or 3 respectively). Preparation-period injury was defined as any reported injury coded 1 or 3; competition-period injury as any reported injury coded 2 or 3. Athletes could therefore contribute a positive outcome to one or both periods.

For each reported injury, athletes also selected the anatomical location (hip, groin, thigh, lower leg, ankle, foot, or head/upper limb/trunk; multi-select) and, only for injuries localized to the lower limb, the injury type (joint sprain, tendon rupture, ligament injury, muscle strain, bone stress injury or stress fracture, tendinopathy, plantar fasciopathy, or other; multi-select, adapted from consensus-recommended categories \citep{bahr2020strobesiis}). These location and type items were not used as covariates or outcomes in the primary Bayesian model; they are reported descriptively (Section~\ref{sec:results-pattern}) to characterize the injury pattern in this cohort, following STROBE-SIIS item SIIS17.1. Injury severity (time-loss duration) was also collected but is outside the scope of the present analysis.

Athletes reporting at least one current-season injury were also presented with four items adapted from the updated OSTRC-H2, covering the injury's impact on sports participation, training volume, sports performance, and symptoms \citep{clarsen2014ostrch,clarsen2020ostrc}. The four ordered responses were scored 0, 8, 17, and 25, then summed to an adapted OSTRC severity score from 0 (no reported impact) to 100 (maximum reported impact). Because these items were administered once retrospectively and only after a current-season injury was reported, the score is used solely as a descriptive athlete-reported impact measure; it is not a prospective weekly OSTRC prevalence or cumulative-burden estimate.

\subsection{Covariates}
\label{sec:covariates}
Adjustment covariates were age (years), competition sex, height (cm), body mass (kg), years of athletics experience, injury in either of the two preceding seasons (binary), period-matched total training-type-frequency scores for each of the five domains and for \emph{palestra}, and indicator variables for the seven discipline categories with non-zero support. These variables were selected a priori based on previously reported risk factors for running-related injury \citep{taunton2003running} and on the requirement, specified in the STROBE-SIIS extension, to account for how athletes were categorized by sport, discipline, and competitive level \citep{bahr2020strobesiis}. Continuous covariates were standardized (zero mean, unit variance) before model fitting.

\subsection{Statistical methods}
\label{sec:statistics}
All inference is Bayesian \citep{gelman2013bda}; no null-hypothesis-testing framework or $p$-values are reported. Results are summarized as posterior means, medians, standard deviations, and 95\% credible intervals (CrI, equal-tailed 2.5th--97.5th percentiles). Every reported interval for a signed association contrast is accompanied by its posterior probability of direction, $P_{\mathrm{dir}}=\max\{P(\Delta>0),P(\Delta<0)\}$ (with one replacing zero as the null for ratios): the posterior probability that the contrast lies on the same side of the null as its median estimate. This is a graded measure of directional certainty, not a $p$-value or a binary significance test, and it must be read together with the effect magnitude and CrI. $P_{\mathrm{dir}}$ is not reported for absolute prevalences, for which positivity is constrained by definition and therefore uninformative. We additionally report posterior probabilities that an association exceeds a pre-specified practically relevant threshold (5 percentage points for differences, a ratio of 1.10) or falls within a region of practical equivalence to no association ($\pm$2 percentage points; ratio 0.95--1.05). Models were fit in Stan \citep{carpenter2017stan} via its R interface.

The adapted OSTRC scores were summarized by mutually exclusive injury-period groups (preparation only, competition only, or both periods) using medians and interquartile ranges. As a descriptive effect-size aid, standardized mean differences (SMD) compared preparation-only and both-period groups with the competition-only group, calculated as the mean-score difference divided by the square root of the average of the two group variances. No injured-versus-uninjured OSTRC comparison was made because the items were not administered to uninjured athletes.

\subsubsection{Primary model}
\label{sec:primary-model}
Each athlete contributed one Bernoulli outcome per period. A shared athlete-level random intercept linked the two outcomes, following a cluster-specific logistic formulation for correlated binary data \citep{neuhaus1991correlated}. For athlete $i$ in period $p$,
\[
Y_{ip}\sim\operatorname{Bernoulli}(\pi_{ip}),
\qquad
\operatorname{logit}(\pi_{ip})=\alpha_p+c_{ip}+a_{ip}+\delta_i+u_i,
\qquad
u_i\sim\mathcal{N}(0,\sigma_u^2).
\]
Here $c_{ip}$ is the carbon-use contribution:
\[
c_{ip}=\begin{cases}
0, & \text{for a non-user},\\
\beta^{\text{any}}_p+f_p(\text{share}_{ip}), & \text{for a carbon user}.
\end{cases}
\]
The remaining terms have direct roles: $\alpha_p$ is the period-specific intercept, $a_{ip}$ combines the athlete-level and period-matched covariates of Section~\ref{sec:covariates}, $\delta_i$ is the partially pooled discipline contribution, and $u_i$ is the athlete random intercept. The function $f_p$ is a period-specific curve over carbon share among users, given a Hilbert-space approximate Gaussian-process (HSGP) prior with a Mat\'ern-5/2 kernel \citep{riutortmayol2023hsgp}, fitted with a common basis and centered at the pooled median user share (0.50):
\[
f_p(\text{share}_{ip}) = \sum_{m=1}^{40} \phi_m(\text{share}_{ip})\, b_{p,m},
\qquad
b_{p,m} \sim \mathcal N\!\left(0,\, S\!\left(\sqrt{\lambda_m};\, \alpha^{\text{gp}}_p, \rho^{\text{gp}}_p\right)\right),
\]
where $\phi_m$ and $\lambda_m$ are the eigenfunctions and eigenvalues of the Laplace operator on the centered, boundary-extended domain $[-0.75, 0.75]$ (a boundary factor of 1.5 times the largest centered observed user share), $S(\cdot)$ is the Mat\'ern-5/2 spectral density, and $\alpha^{\text{gp}}_p$ and $\rho^{\text{gp}}_p$ are period-specific marginal-standard-deviation and length-scale hyperparameters estimated from the data. This construction lets the posterior determine how much curvature the competition- and preparation-period dose--response relationships each support directly from the data, through $\alpha^{\text{gp}}_p$ and $\rho^{\text{gp}}_p$. 40 basis functions were confirmed sufficient by checking that the spectral density at the highest included frequency was negligible relative to the lowest across both hyperparameters' posterior distributions. Consequently, $\beta^{\text{any}}_p$ compares a median-share user with a non-user on the conditional log-odds scale, while $f_p$ describes dose among users. The athlete and discipline terms, and the HSGP basis coefficients $b_{p,m}$, were parameterized non-centered, with discipline coefficients partially pooled through a shared scale $\tau_{\text{disc}}$. Priors were regularizing and weakly informative on the standardized log-odds scale: $\alpha_p\sim\mathcal{N}(\operatorname{logit}(0.4),1.5^2)$, $\beta^{\text{any}}_p\sim\mathcal N(0,0.7^2)$, covariate coefficients $\mathcal{N}(0, 0.5^2)$, $\tau_{\text{disc}} \sim \text{Half-N}(0, 0.5^2)$, $\sigma_u \sim \text{Half-N}(0, 1^2)$, $\alpha^{\text{gp}}_p \sim \text{Half-N}(0, 0.7^2)$, and $\rho^{\text{gp}}_p \sim \text{Inv-Gamma}(2.94, 0.80)$, the last chosen so that its 1st and 99th prior percentiles fell at one-tenth and twice the observed centered-share range, discouraging length scales too short for the data to resolve while still permitting substantial curvature. A matched baseline model retained the identical outcome, covariate, discipline, and athlete-random-intercept structure but removed both carbon-use terms, providing a no-exposure comparator for predictive-validity assessment (Section~\ref{sec:validation}).

Figure~\ref{fig:primary-model-graph} summarizes the primary model's observed inputs, parameter hierarchy, and repeated athlete-period outcome structure.
\begin{figure}[htbp]
\centering
\resizebox{\textwidth}{!}{%
\begin{tikzpicture}[
  >=Latex,
  every node/.style={font=\normalsize},
  latent/.style={draw, ellipse, fill=blue!7, minimum height=10mm, align=center, inner sep=2pt},
  data/.style={draw, rounded corners=2pt, fill=black!8, minimum height=10mm, align=center, inner sep=2pt},
  term/.style={draw, rounded corners=2pt, fill=green!8, minimum height=10mm, align=center, inner sep=2pt},
  outcome/.style={draw, ellipse, fill=black!18, minimum height=12mm, align=center, inner sep=2pt},
  summary/.style={draw, dashed, rounded corners=2pt, fill=orange!8, minimum height=10mm, align=center, inner sep=2pt},
  plate/.style={draw, rounded corners=4pt, inner sep=6mm},
  arrow/.style={->, thick}
]


  \node[latent]                  (alpha)        at (-1.0, 0)     {$\alpha_p$\\$\mathcal N(\operatorname{logit}(0.4),1.5^2)$};
  \node[latent]                  (exposurecoef) at (4.7, 0)      {$\beta^{\mathrm{any}}_p,\boldsymbol\beta^{\mathrm{share}}_p$\\$\mathcal N(0,0.7^2)$};
  \node[latent]                  (commoncoef)   at (10.0, 0)     {$\boldsymbol\beta$\\$\mathcal N(0,0.5^2)$};
  \node[latent]                  (discscale)    at (13.5, 0)     {$\tau_{\mathrm{disc}}$\\$\operatorname{Half\text{-}N}(0,0.5^2)$};
  \node[latent]                  (athscale)     at (17.5, 0)     {$\sigma_u$\\$\operatorname{Half\text{-}N}(0,1^2)$};


  \node[data]                    (period)       at (-1.0, -2.3)  {period $p$\\preparation / competition};
  \node[data]                    (adoptioninput) at (3.0, -2.3)  {any carbon use\\$A_{ip}=\mathbb 1(s_{ip}>0)$};
  \node[data]                    (doseinput)     at (6.4, -2.3)  {carbon share\\$s_{ip}$; basis $\mathbf B(s_{ip})$};
  \node[data]                    (covariates)   at (10.0, -2.3)  {covariates\\$\mathbf X_{ip}$};
  \node[latent]                  (disccoef)     at (13.5, -2.3)  {$\delta_d$, $d=1,\ldots,7$\\$\mathcal N(0,\tau_{\mathrm{disc}}^2)$};
  \node[latent]                  (athleteeffect) at (17.5, -2.3) {$u_i$, $i=1,\ldots,352$\\$\mathcal N(0,\sigma_u^2)$};


  \node[term]                    (adoptionterm) at (3.0, -4.6)   {adoption\\$A_{ip}\beta^{\mathrm{any}}_p$};
  \node[term]                    (doseterm)     at (6.4, -4.6)   {dose among users\\$\mathbf B(s_{ip})\boldsymbol\beta^{\mathrm{share}}_p$};


  \node[term]                    (interceptterm)  at (-1.0, -6.6) {$\alpha_p$};
  \node[term]                    (exposureterm)   at (4.7, -6.6)  {carbon contribution\\$c_{ip}=c^{\mathrm{adoption}}_{ip}+c^{\mathrm{dose}}_{ip}$};
  \node[term]                    (covariateterm)  at (10.0, -6.6) {$\mathbf X_{ip}\boldsymbol\beta$};
  \node[term]                    (discterm)       at (13.5, -6.6) {$\mathbf Z_i\boldsymbol\delta$};
  \node[term]                    (athleteterm)    at (17.5, -6.6) {$u_i$};


  \node[term, minimum width=48mm] (periodassembly) at (9.0, -9.1)
    {evaluate the linear predictor\\separately for each period};

  \node[term, minimum width=38mm] (etaprep) at (5.8, -11.2)
    {$\eta_{i,\mathrm{prep}}=\operatorname{logit}(\pi_{i,\mathrm{prep}})$};
  \node[term, minimum width=38mm] (etacomp) at (12.2, -11.2)
    {$\eta_{i,\mathrm{comp}}=\operatorname{logit}(\pi_{i,\mathrm{comp}})$};


  \node[outcome] (prepoutcome) at (5.8, -13.4) {$Y_{i,\mathrm{prep}}$\\$\operatorname{Bernoulli}(\pi_{i,\mathrm{prep}})$};
  \node[outcome] (compoutcome) at (12.2, -13.4) {$Y_{i,\mathrm{comp}}$\\$\operatorname{Bernoulli}(\pi_{i,\mathrm{comp}})$};


  \node[summary] (anyoutcome) at (9.0, -16.2) {descriptive any-period injury\\$Y_i^{\mathrm{any}}=\max(Y_{i,\mathrm{prep}},Y_{i,\mathrm{comp}})$};

  \begin{pgfonlayer}{background}

    \draw[arrow] (alpha) -- (interceptterm);
    \draw[arrow] (period) -- (interceptterm);


    \draw[arrow] (exposurecoef.south west) to[out=-100,in=90] (adoptionterm.north);
    \draw[arrow] (exposurecoef.south east) to[out=-80,in=90] (doseterm.north);
    \draw[arrow] (adoptioninput) -- (adoptionterm);
    \draw[arrow] (doseinput) -- (doseterm);
    \draw[arrow] (adoptionterm) -- (exposureterm.north west);
    \draw[arrow] (doseterm) -- (exposureterm.north east);


    \draw[arrow] (commoncoef) -- (covariateterm);
    \draw[arrow] (covariates) -- (covariateterm);


    \draw[arrow] (discscale) -- (disccoef);
    \draw[arrow] (disccoef) -- (discterm);


    \draw[arrow] (athscale) -- (athleteeffect);
    \draw[arrow] (athleteeffect) -- (athleteterm);


    \draw[arrow] (interceptterm.south) to[out=-90,in=160] (periodassembly.west);
    \draw[arrow] (exposureterm.south) to[out=-90,in=150] (periodassembly.north west);
    \draw[arrow] (covariateterm) -- (periodassembly.north);
    \draw[arrow] (discterm.south) to[out=-90,in=30] (periodassembly.north east);
    \draw[arrow] (athleteterm.south) to[out=-90,in=20] (periodassembly.east);


    \draw[arrow] (periodassembly.south west) to[out=-100,in=90] (etaprep.north);
    \draw[arrow] (periodassembly.south east) to[out=-80,in=90] (etacomp.north);

    \draw[arrow] (etaprep) -- (prepoutcome);
    \draw[arrow] (etacomp) -- (compoutcome);


    \draw[arrow] (prepoutcome.south) to[out=-90,in=140] (anyoutcome.north west);
    \draw[arrow] (compoutcome.south) to[out=-90,in=40] (anyoutcome.north east);
  \end{pgfonlayer}

  \node[plate, fit=(periodassembly)(etaprep)(etacomp)(prepoutcome)(compoutcome)] (modelplate) {};
  \node[anchor=west, align=left] at ([xshift=3mm]modelplate.east) {$352$ athletes\\$\times\,2$ periods};
\end{tikzpicture}%
}
\caption{Graphical representation of the primary Bayesian multilevel model. Grey boxes are observed inputs, blue ellipses are estimated parameters, green boxes are deterministic contributions to the linear predictor, and shaded ellipses are the two observed period-specific injury outcomes, jointly enclosed in the solid plate (352 athletes $\times$ 2 periods). The carbon contribution $c_{ip}$ combines adoption ($A_{ip}$, any carbon use) and dose ($\mathbf B(s_{ip})$, the centered three-degree-of-freedom natural-spline basis for carbon share among users). $\mathbf Z_i$ denotes the seven discipline indicators underlying the partially pooled coefficients $\delta_d$, and $u_i$ the shared athlete random intercept. The common model structure is evaluated separately as $\eta_{i,\mathrm{prep}}$ and $\eta_{i,\mathrm{comp}}$; each predictor points only to its corresponding outcome. The dashed box is the derived any-period injury indicator reported descriptively, outside the probabilistic-model plate, and is not an additional likelihood term.}
\label{fig:primary-model-graph}
\end{figure}


\subsubsection{Estimands and marginal standardization}
\label{sec:estimands}
Posterior draws were used to compute marginal (population-standardized) adjusted prevalences by averaging model predictions over the observed empirical distribution of covariates and discipline indicators within each period, and by integrating the athlete random intercept using 15-point Gauss--Hermite quadrature rather than conditioning on a fixed value of $u_i$. For each period, adjusted prevalence was computed at four exposure scenarios: non-user, and user at the empirical 25th (P25), 50th (P50), and 75th (P75) percentile of the observed among-user share distribution for that period. These are ranks in the observed exposure distribution, not fixed share values: P25, P50, and P75 corresponded to 28.6\%, 42.1\%, and 57.1\% share in preparation and to 36.0\%, 50.0\%, and 71.4\% share in competition. Two primary contrasts were derived: the \emph{adoption} contrast (P50 user vs.\ non-user; prevalence difference and ratio) and the \emph{dose} contrast (P75 vs.\ P25 among users; prevalence difference and ratio).

To provide a directly interpretable absolute-risk presentation without fitting an additional model, the same posterior draws were also evaluated for one fixed reference profile. The profile was a male speed athlete aged 21 years, 176~cm tall, weighing 64~kg, with 10 years of athletics experience and an injury in either of the two preceding seasons. The period-matched training-frequency scores were fixed at their pooled cohort medians: long endurance 1, medium endurance 1, lactate 2, maximum velocity 2, technique 2, and gym 2; these are the survey's ordinal frequency scores, not exact counts of unique sessions. The athlete random intercept was fixed at zero (the median of its modeled distribution). Thus, ``median profile'' denotes medians for continuous and ordinal covariates combined with the modal observed sex and discipline profile, not a uniquely defined median athlete. Adoption was summarized as the risk difference between a non-user and a user at 50\% carbon share. Among users, the continuous HSGP risk curve was evaluated over the share range supported in both periods, with preparation and competition predictions calculated for this identical profile.

For display only, each continuous curve was read at three exact share values: 25\%, 50\%, and 75\%. These values are not P25, P50, and P75 of the sample, do not define exposure groups, and are not proposed clinical thresholds. They were selected after model fitting as round, evenly spaced interpretive anchors within the observed support of both periods; 50\% is also the pooled median user share and the centering point of the HSGP curve, while 25\% and 75\% are symmetric around it. A change from 25\% to 50\%, or from 50\% to 75\%, therefore means changing the carbon component from one-quarter to one-half, or from one-half to three-quarters, of the same fixed total reported training-frequency score for the same reference profile. Both changes are derived from one continuous fitted function, not from separate models or comparisons between athlete subgroups. These are conditional, profile-specific model predictions intended to aid interpretation of the fitted association, not validated individual prognostic estimates.

\subsubsection{Matched comparison of relative and absolute exposure}
\label{sec:exposure-scale-comparison}
To examine whether the observed association was specific to relative allocation or instead tracked absolute carbon use, two matched multilevel models represented two substantive exposure hypotheses. The \emph{relative-use hypothesis} used carbon share and asks whether injury prevalence changes as a greater proportion of a given training-frequency total is allocated to carbon-plated spikes. The \emph{absolute-use hypothesis} used the absolute carbon-frequency score and asks whether injury prevalence changes with more reported carbon use after adjustment for the same total training-frequency covariates. Both exposures were standardized among users across both periods (mean 0, standard deviation 1) and entered linearly with period-specific coefficients; each model also retained the period-specific adoption term, identical covariates, discipline effects, priors, and athlete random-intercept structure of the primary model. One standard deviation represented 0.241 in share (24.1 percentage points) or 3.36 units of the summed ordinal carbon-frequency score. The latter is a change in the questionnaire-derived score, not 3.36 unique training sessions, because frequency categories were top-coded and training domains could overlap within a session. The two exposures were fitted in separate models because absolute carbon frequency contributes directly to both the share numerator and the adjusted total-frequency scores; simultaneous inclusion would make the coefficients difficult to identify and interpret.

For a common conditional estimand, posterior injury risk was evaluated from 0.5 standard deviations below to 0.5 standard deviations above the pooled user mean. On the observed scales, this one-standard-deviation increase was from 37.3\% to 61.4\% carbon share or from 2.83 to 6.19 units of the summed ordinal carbon-frequency score. Continuous covariates were fixed at their cohort medians, discipline at the most prevalent observed indicator, and the athlete random intercept at zero. The comparison was designed to place the two scales on equal footing, not to treat absence of an association in one model as proof of no effect or to select a causal exposure mechanism. This matched exposure-scale comparison was added to clarify interpretation of the primary share analysis and was not part of the prospectively specified sensitivity set.

\subsubsection{Validation and sensitivity analyses}
\label{sec:validation}
Predictive comparison between the primary (exposure) and baseline (no-exposure) models used Pareto-smoothed importance-sampling leave-one-out cross-validation (PSIS-LOO) \citep{vehtari2017loo}, computed on the athlete-grouped joint log-likelihood (both period records per athlete evaluated jointly, with the athlete random intercept marginalized by Gauss--Hermite quadrature before leave-one-athlete-out prediction), reported as the expected log predictive density (ELPD) difference with its standard error. Posterior predictive checks compared observed and model-replicated counts of injured athletes per period. Two pre-specified sensitivity analyses tested structural assumptions of the primary specification: (1) a linear carbon-share term in place of the HSGP curve; (2) wider hyperpriors (HSGP marginal standard deviation $\alpha^{\text{gp}}_p \sim \text{Half-N}(0, 1.2^2)$, covariate coefficients $\mathcal{N}(0, 1.0^2)$, length-scale prior unchanged). Sensitivity models were compared with the primary model as robustness checks on the same estimands, not as competing substantive hypotheses.

\subsection{Population reweighting}
\label{sec:weighting}
Because the two-phase recruitment (Section~\ref{sec:participants}) could not guarantee equal selection probabilities across age-by-sex strata, a poststratification step \citep{gelman2007poststrat} reweighted the analysis sample to FIDAL national population estimates for six age-by-sex strata (Under 20, Under 23, Senior $\times$ Men, Women; Appendix Table~\ref{tab:weight-diagnostics}). Stratum model weights were the ratio of the population share to the observed sample share, further regularized to avoid extreme weights. The Kish effective sample size after weighting was 340.5 of 352 (97\%), indicating a modest efficiency loss from reweighting. Weighted estimands were computed by combining the athlete-level posterior predictions of Section~\ref{sec:estimands} with model weights before marginalization, and are reported alongside the unweighted (sample-representative) estimands as a sensitivity check on generalizability to the national competitive population.

\subsection{Subgroup analyses}
\label{sec:subgroups}
As a pre-specified descriptive and exploratory extension, the primary adjusted estimands were recomputed conditional on sex (men, women) and on age group (Under 20, Under 23, Senior), holding the same model structure and covariate adjustment, to assess whether the direction and approximate magnitude of the competition-period dose association were consistent across these strata. These subgroup estimates are not powered as confirmatory tests of effect-modification and are reported as descriptive, hypothesis-generating comparisons.

\subsection{Sources of bias}
\label{sec:bias}
The survey instrument underwent content-validity and face-validity review by athletics-medicine content experts and prospective respondents before deployment, but no formal test-retest reliability or criterion-validity assessment against objectively measured training data (for example, coach-logged sessions or shoe-embedded wear sensors) was conducted; this is treated as a limitation rather than a validated measurement property (Section~\ref{sec:discussion-limitations}). Both exposure and outcome were self-reported for the same recall period by the same respondent, which admits recall bias and the possibility that athletes' beliefs about carbon-plate risk influenced how they reported either quantity. The survey was administered near the end of the competitive season, so the competition-period recall window was recent while the preparation-period window reached further back into the season; this asymmetry is a plausible source of differential, period-specific recall error, and any resulting attenuation in preparation-period exposure and outcome reporting would bias that period's association toward the null independently of the underlying biomechanics. This offers one plausible explanation, among others, for why the preparation-period association was weaker and less certain than the competition-period one, and is a reason to treat the period contrast itself, not only the individual period estimates, with some caution. Reporting was not restricted to medically trained staff; athletes reported their own injuries and training without independent clinical verification, following the broad, non-medical-attention-dependent injury definition of Section~\ref{sec:outcome}, which was chosen specifically to reduce under-ascertainment of injuries that did not prompt a clinical visit at the cost of relying on athlete-reported symptom recognition rather than diagnosis.

\subsection{Missing data and other analyses}
\label{sec:missing}
The analysis used a complete-case cohort by design: eligibility required a fully completed survey, and the three further exclusions in Section~\ref{sec:participants} addressed incomplete injury timing or a structurally undefined carbon-share denominator. No further imputation was performed.

\subsection{Software}
\label{sec:software}
The fitted models used a Bernoulli likelihood with a logit link. Posterior sampling used Stan's No-U-Turn Sampler (NUTS), a Hamiltonian Monte Carlo algorithm, through the \texttt{rstan} interface. For each model, four independent chains each ran for 2000 iterations; the first 1000 were used for warm-up and adaptation and discarded, leaving 1000 retained draws per chain and 4000 retained draws in total to approximate the posterior distribution. The target acceptance probability was set to \texttt{adapt\_delta}=0.98 and the maximum tree depth to 12. Runs used base seed 20260719 with deterministic model-specific offsets. Analyses used the \texttt{loo} package for PSIS-LOO cross-validation; the HSGP basis, spectral density, and length-scale prior calibration were implemented directly in Stan and base R, with no additional package dependency. Full software versions are reported in Appendix~\ref{sec:session-info}.
